% =========================================================================
% Capítulo 1: Cinemática
% Tema 04: Rodadura sin deslizar
% Problema 04: Sistema de dos discos rodantes
% =========================================================================
\subsubsection{Sistema de dos discos rodantes}

Un disco circular de radio $R$ rueda sin resbalar a lo largo de una superficie horizontal fija. Un segundo disco circular de radio $r$ rueda sin resbalar sobre el primer disco. Usando las variables $\theta$ y $\phi$ para describir el movimiento de los discos, determine las siguientes cantidades como las ve un observador en un marco de referencia inercial:

\begin{enumerate}
    \item Velocidad angular y aceleración angular de cada disco.
    \item Velocidad y aceleración del centro de cada disco ($O$ y $C$).
\end{enumerate}

\begin{figure}[!ht]
    \centering
    \begin{tikzpicture}[>=Stealth, scale=1.5]
        % --- Parámetros Geométricos ---
        \def\Rbig{2.0}    % Radio del disco inferior (R)
        \def\Rsmall{1.0}  % Radio del disco superior (r)
        \def\ThetaPos{30} % Ángulo de posición del centro C respecto a la vertical
        \def\PhiPos{40}   % Ángulo interno phi del disco superior

        % --- 1. Suelo Fijo ---
        % Línea base
        \draw[thick] (-3, 0) -- (4, 0);
        % Patrón de sombreado para indicar superficie fija
        \fill[pattern=north east lines, pattern color=gray!50] (-3, -0.3) rectangle (4, 0);

        % --- 2. Disco Inferior (Centro O) ---
        \coordinate (O) at (0, \Rbig);
        % Sombreado radial para efecto 3D
        \shadedraw[inner color=white, outer color=gray!40, thick] (O) circle (\Rbig);
        \filldraw (O) circle (1.5pt) node[left, xshift=-2pt] {$O$};

        % --- 3. Disco Superior (Centro C) ---
        % Calculamos la posición de C basándonos en el ángulo ThetaPos
        \coordinate (C) at ($(O) + ({90-\ThetaPos}:{\Rbig+\Rsmall})$);
        \shadedraw[inner color=white, outer color=gray!40, thick] (C) circle (\Rsmall);
        \filldraw (C) circle (1.5pt) node[above, yshift=2pt] {$C$};

        % --- 4. Punto de Contacto P ---
        \coordinate (P) at ($(O) + ({90-\ThetaPos}:\Rbig)$);
        \filldraw (P) circle (1.2pt) node[right, xshift=3pt] {$P$};

        % --- 5. Líneas de Referencia y Radios ---
        % Línea vertical desde O
        \draw[dashed, gray] (O) -- ++(0, {\Rbig+\Rsmall+1.0});
        % Línea conectando centros OC
        \draw[dashed, gray] (O) -- (C);
        
        % Etiqueta Radio R (disco inferior)
        \draw[->, thin] (O) -- ++(-30:\Rbig) node[midway, below, scale=0.9] {$R$};
        % Etiqueta Radio r (disco superior)
        \draw[->, thin] (C) -- ++({90-\ThetaPos+50}:\Rsmall) node[midway, above right, scale=0.9] {$r$};

        % --- 6. Ángulos ---
        % Ángulo Theta (posición de C respecto a la vertical)
        \draw[->, thick] ($(O)+(0, 1.5)$) arc (90:{90-\ThetaPos}:1.5);
        \node at ($(O)+({90-0.5*\ThetaPos}:1.7)$) {$\phi$};

        % Ángulo Phi (interno en el disco superior)
        % Dibujamos una línea radial de referencia en el disco superior
        \coordinate (CRef) at ($(C) + ({90-\ThetaPos-\PhiPos}:\Rsmall)$);
        \draw[dashed, gray!80] (C) -- (CRef);
        % Arco para phi desde la línea CP hasta CRef
        %\draw[->] ($(C)!0.5!(P)$) arc ({90-\ThetaPos}:{90-\ThetaPos-\PhiPos}:0.5);
        %\node at ($(C)+({90-\ThetaPos-0.5*\PhiPos}:0.7)$) {$ $};

        % --- 7. Flechas de Sentido de Rotación (Decorativas) ---
        % Disco inferior
        \draw[<-, thick, bend right] ($(O)+(-0.6, \Rbig+0.5)$) to node[midway, above] {$\theta$} ($(O)+(-1.5, \Rbig)$);
        % Disco superior
        \draw[<-, thick, bend left] ($(C)+(\Rsmall+0.3, 0.3)$) to ($(C)+(\Rsmall+0.3, -0.3)$);

        % --- 8. Indicación "No Slip" ---
        \node[right] at (3, \Rbig) {Sin deslizamiento};
        % Línea al contacto inferior
        \draw[thin] (3, \Rbig) -- (0,0);
        % Línea al contacto intermedio P
        \draw[thin] (3, \Rbig) -- (P);

    \end{tikzpicture}
    \caption{Sistema de dos discos apilados rodando sin deslizamiento.}
    \label{fig:dos_discos}
\end{figure}